\section{10. Limitations}

\subsection{10.1 BAM-UMoE is Not a Compression Method}

\textbf{Critical.} BAM-UMoE \textbf{increases} parameter count by ~350\%. 
It should never be described as "compression" or "efficiency". Any paper 
using BAM-UMoE must disclose the parameter increase.

\subsection{10.2 Parameter Overhead}

The MoE branch adds $O(E \cdot d^2)$ parameters per layer:
\begin{itemize}
\item $d = 768$, $E = 8$: ~37M params per layer.
\end{itemize}
\begin{itemize}
\item 6-layer DistilBERT: +226M params.
\end{itemize}

\subsection{10.3 Compute Overhead}

The MoE branch adds $O(L d^2)$ FLOPs per layer, comparable to base 
attention. Total compute roughly doubles.

\subsection{10.4 VRAM Overhead}

The MoE branch's expert weights and activations are stored in VRAM, 
increasing peak memory by ~{{VRAM_INC}}×.

\subsection{10.5 Inference Latency}

Two matmuls + router overhead = ~{{LATENCY_RATIO}}× slower.

\subsection{10.6 Task Coverage}

Only SST-2 evaluated. Generation, QA, NLI untested.

\subsection{10.7 Model Coverage}

Three models tested (66.9M–124.6M). Larger models untested.

\subsection{10.8 Hardware}

Only RTX 4070 Laptop. Scaling to datacenter GPUs unknown.

\subsection{10.9 Load Balancing}

Auxiliary loss weight $\lambda_\text{aux} = 0.01$ chosen 
heuristically. Not tuned extensively.

\subsection{10.10 Single Seed}

Standard error ~1.69\%. Differences below 3.4\% may be noise.

\subsection{10.11 Cross-Attention Scope}

MoE is not applied to cross-attention. This is a design choice but 
limits applicability to encoder-decoder tasks.

\subsection{10.12 Code Quality}

Research-grade:
\begin{itemize}
\item Broad exception handling
\end{itemize}
\begin{itemize}
\item No unit tests for MoE routing
\end{itemize}
\begin{itemize}
\item Numerical stability could be improved
\end{itemize}

\subsection{10.13 Ethical Considerations}

Dual-use. See Appendix E.